\documentclass[10pt]{article}
\providecommand{\FIGDIR}{fig}
\newcommand{\DIGESTDATE}{24 October 2026}
\newcommand{\PERIODSTART}{18 October 2026}
\newcommand{\PERIODEND}{24 October 2026}
\input{preamble}

\begin{document}

% ==================================================================== MASTHEAD
\thispagestyle{empty}
\begin{tikzpicture}[remember picture, overlay]
  \fill[Deep] (current page.north west) rectangle
      ($(current page.north east)+(0,-67mm)$);
  \fill[Teal] ($(current page.north west)+(0,-67mm)$) rectangle
      ($(current page.north east)+(0,-68.4mm)$);
  \foreach \x/\y/\r in {22/12/0.9, 41/26/1.5, 63/9/0.7, 88/31/1.1, 112/17/1.9,
                        138/29/0.8, 159/13/1.3, 178/54/1.0, 196/21/1.6,
                        30/44/0.6, 74/57/1.2, 125/48/0.9, 168/6/0.8, 191/59/0.7,
                        52/18/1.0, 100/60/0.6, 148/20/1.1, 66/49/1.4, 205/45/0.8}
    \fill[Sky, opacity=0.30] ($(current page.north west)+(\x mm,-\y mm)$) circle (\r mm);
\end{tikzpicture}

\vspace*{4mm}
{\sffamily\fontsize{7.6}{9}\selectfont\color{Sky}%
 WEEKLY ROLLUP\;\textbullet\;PhD RESEARCH WATCH\par}
\vspace{2.2mm}
{\sffamily\bfseries\fontsize{27}{29}\selectfont\color{white}%
 Particulate Matter\par}
{\sffamily\fontsize{27}{29}\selectfont\color{Sky}%
 Week in Review\par}
\vspace{5.0mm}
{\sffamily\fontsize{8.4}{10}\selectfont\color{white}%
 \mbox{Week ending \DIGESTDATE}\;\;\textcolor{Teal}{\textbar}\;\;\mbox{\PERIODSTART{} to \PERIODEND}\;\;
 \textcolor{Teal}{\textbar}\;\;\mbox{119 records pooled from 7 daily issues}\par}
\vspace{18mm}

% -------------------------------------------------------------- metric strip
\noindent
\begin{minipage}[t]{0.190\linewidth}\metric{119}{RECORDS POOLED\\(119 UNIQUE DOIs)}{Deep}\end{minipage}\hfill
\begin{minipage}[t]{0.190\linewidth}\metric{39\%}{MEASUREMENT SIDE\\(46 RECORDS)}{Teal}\end{minipage}\hfill
\begin{minipage}[t]{0.190\linewidth}\metric{10}{SUBTOPICS\\ACTIVE}{Sky}\end{minipage}\hfill
\begin{minipage}[t]{0.190\linewidth}\metric{30}{POOLED EFFECT\\ESTIMATES (30 WITH A CI)}{Amber}\end{minipage}\hfill
\begin{minipage}[t]{0.190\linewidth}\metric{25}{ABSTRACT-FREE\\(21\% OF WEEK)}{Clay}\end{minipage}

\vspace{6mm}

% -------------------------------------------------------------- week in one box
\section*{\sffamily\bfseries\color{Deep}The week in one paragraph}
\vspace{-1.5mm}{\color{Teal}\rule{\linewidth}{1.1pt}}\vspace{2.5mm}

\begin{keybox}
{\sffamily\bfseries\small\color{Deep}Volume steady at 119, tier A at a series high of 18\,\%, and the week's
epidemiology split PM$_{2.5}$ into its parts: sulfate, nitrate and ammonium carried the cardiorenal, cognitive and
atrial-fibrillation signals, and elemental carbon carried ADHD. On the measurement side, a meta-analysis put a
number on the optical-monitor bias every calibration inherits.}\\[1.4mm]
\footnotesize
\textbf{Scale.} \textbf{119 records} over seven issues on daily totals of 13, 18, 19, 18, 17, \textbf{24} and
\textbf{10}. Against recent rollups: 119 after 123, 87, 125, 162, 148 and 126. \textbf{All ten subtopics are
populated} for the fifth week running.

\textbf{The measurement side held at 39\,\%, but its make-up flipped.} \textit{Exposure assessment \&
modelling} 26 $\rightarrow$ \textbf{31} (forecasting emulators, CTM and emission-factor work); \textit{Sensing}
22 $\rightarrow$ \textbf{15}. Share series 44, 33, 42, 44, 39, \textbf{39\,\%}.

\textbf{Cardiovascular \& metabolic led the health side's growth}: 7 $\rightarrow$ \textbf{12};
\textit{Respiratory} 8 $\rightarrow$ \textbf{11}; \textit{Burden} 18 $\rightarrow$ 14 and \textit{Mechanistic
toxicology} 15 $\rightarrow$ 10 contracted.

\textbf{Estimates: 30, every one with an interval} (fifth complete week at 100\,\%). One per 4.0 records.
27 sit on ratio scales (HR 16, RR 6, OR 5) and 3 on difference scales (FEV$_1$ in mL, HAZ, and a log-ratio of
optical to gravimetric mass).

\textbf{Abstract gap 25 of 119 (21\,\%)}, unchanged from last week. \textbf{Tier mix 21 A, 43 B, 55 C} ---
tier A at \textbf{18\,\%}, up from 12\,\%.
\end{keybox}

\clearpage

% -------------------------------------------------------------- trend
\section*{\sffamily\bfseries\color{Deep}Subtopic trend and inflection}
\vspace{-1.5mm}{\color{Teal}\rule{\linewidth}{1.1pt}}\vspace{2.5mm}

\noindent\includegraphics[width=\linewidth]{\FIGDIR/w1_subtopic_trend.png}
\figcap{Per-issue subtopic counts across the week from \texttt{state/metrics.csv}, seven issues. Daily
totals 13, 18, 19, 18, 17, 24 and 10. 18 October was built on the 19 October run, 19--22 October on their
next-day runs; the scheduled runs for 23--29 October did not complete, so 23 and 24 October and this rollup were
built on the \textbf{30 October recovery run}.}

\footnotesize
\textbf{Expanding.} \textit{Exposure} 26 $\rightarrow$ \textbf{31} is the largest absolute gain and is
forecasting-heavy: a physics--AI coupled CTM forecaster over 338 cities, a km-scale generative global PM$_{2.5}$
forecaster (Foehn), and a GEOS-CF emulator whose skill collapsed under blocked spatial cross-validation (RF on
latitude/longitude alone R$^2$ 0.91; $r$ 0.11 on an unseen sector). \textit{Cardiovascular \& metabolic}
7 $\rightarrow$ \textbf{12} carries the week's largest cohorts.

\vspace{1.2mm}
\textbf{Contracting.} \textit{Sensing} 22 $\rightarrow$ \textbf{15}, \textit{Mechanistic toxicology}
15 $\rightarrow$ 10, \textit{Burden} 18 $\rightarrow$ 14. \textit{Neuro} flat at 4, but two of the four are
tier A or carry components (Danish ADHD; CLHLS cognition).

\vspace{1.2mm}
\textbf{The inflection worth naming.} \textbf{Composition displaced mass as the exposure that carries the
effect.} In one week: incident eGFR decline in 307,639 adults at HR \textbf{1.273} (sulfate), \textbf{1.304}
(nitrate) and \textbf{1.280} (ammonium); cognitive disorders in CLHLS older adults at HR 1.030--1.033 for the
three secondary inorganic ions against \textbf{1.005} for PM$_{2.5}$ mass; incident atrial fibrillation in 26.3
million Medicare beneficiaries at HR \textbf{1.056} per unit sulfate-source PM$_{2.5}$ vs \textbf{1.029} for
vehicle PM$_{2.5}$; and a second-year-of-life ADHD window that follows \textbf{elemental carbon} (HR 1.24) and is
unchanged after PM$_{2.5}$ adjustment. Every component exposure here is modelled (CTM, satellite--ML or source
apportionment), not measured. Last week this channel was the smoke carbonaceous fraction; this week it is
secondary inorganics. \textbf{The test for next week}: a component-resolved estimate built on measured
speciation, or a network-scale absorption (BC) channel used as the exposure.

\vspace{2mm}
\noindent\includegraphics[width=\linewidth]{\FIGDIR/f2a_architecture.png}
\figcap{Study architecture over the pooled week. \textit{Metadata only} 25 leads, then
\textit{Modelling / inventory} 19, \textit{Review / synthesis} 12, \textit{Measurement campaign} 12 and
\textit{Observational -- cohort} 10. \textbf{Three interventional records}: the HEPA
vs sham crossover (cytokines, visceral adiposity, 18 October), a traffic-noise/air-pollution randomised
crossover (lung function), and the HAPIN LPG follow-up to age 5 (null for HAZ).}

\clearpage
% -------------------------------------------------------------- pooled analytics
\section*{\sffamily\bfseries\color{Deep}Pooled analytics}
\vspace{-1.5mm}{\color{Teal}\rule{\linewidth}{1.1pt}}\vspace{3mm}

\noindent\includegraphics[width=\linewidth]{\FIGDIR/f1_subtopics.png}
\figcap{Pooled subtopic distribution, all ten bands populated: \textit{Exposure} 31 and \textit{Sensing} 15
put the measurement side at \textbf{46 of 119 (39\,\%)}. On the health side \textit{Burden} 14 and
\textit{Cardiovascular} 12 lead, then \textit{Respiratory} and \textit{Occupational \& indoor} 11 each,
\textit{Mechanistic toxicology} 10; \textit{Neuro} is smallest at 4.}

\vspace{2mm}
\noindent\includegraphics[width=\linewidth,height=0.80\textheight,keepaspectratio]{\FIGDIR/f2b_endpoint.png}
\figcap{Health endpoint over the pooled week, canonicalised. ``No health endpoint'' pools the metadata-only
records with instrument, forecasting, emissions and modelling work. Among human outcomes, renal, cognitive,
cardiac-rhythm and mortality endpoints lead; antimicrobial resistance and cervical-cancer recurrence are new
endpoint classes for the watch.}

\vspace{2mm}
\noindent\includegraphics[width=\linewidth]{\FIGDIR/f3_metric_geography.png}
\figcap{Left --- particle metric: \textbf{PM$_{2.5}$ alone 51 of 119}, unspeciated PM or emissions 32,
joint PM$_{2.5}$ + PM$_{10}$ 14, optical properties 6, size-resolved 5, PM$_{10}$ alone 4, composition 2. The
component-resolved epidemiology of the week is coded under PM$_{2.5}$, so the composition bar understates it. Right --- geography:
\textit{Global / multi-region} 34 (reviews, laboratory work, the six- and 102-country records), \textbf{China 32},
Europe 14, North America 12, Southeast Asia 7, South Asia 5, MENA 4, East Asia (ex-China) 4.}

\vspace{2mm}
\noindent\includegraphics[width=\linewidth]{\FIGDIR/f6_lifecourse.png}
\figcap{Life-course windows across the pooled week; counts non-exclusive. \textbf{Working age 22}, older
adults 10, preconception / in utero 7, infancy \& early childhood 7, \textbf{adolescence 1} --- the adolescent
window fell back to a single record after last week's 4.}

\clearpage

\noindent\includegraphics[width=\linewidth]{\FIGDIR/f5_forest.png}
\figcap{The \textbf{27} ratio-scale estimates, every one with a 95\,\% CI. \textbf{Not mutually comparable and
not meta-analysed}: increments include per 1\,$\mu$g\,m$^{-3}$, per IQR, quartile and binary contrasts, a
unitless trajectory index and an MR instrument, on HR, OR and RR scales, so no pooled estimate or $I^2$ is
reported. Within-family heterogeneity is low where it can be read: the three cardiorenal ion HRs (1.273--1.304)
and the three CLHLS ion HRs (1.020--1.033) overlap almost entirely, consistent with collinear modelled
components rather than independent effects. \textbf{Two intervals cross unity}: \textit{S.\,enterica}
tetracycline resistance RR 1.65 (0.98--2.76) and SO$_2$ stroke recurrence HR 1.19 (0.94--1.50). The Beijing
per-1\,$\mu$g\,m$^{-3}$ mortality HR (1.053) is an order of magnitude above pooled long-term estimates. Nested
families share cases: cardiorenal 3, CLHLS 4, Medicare AF 2, cervical cancer 2, Danish ADHD 2, Beijing 2,
Guangzhou SGA 3, Medicare smoke waves 2, antimicrobial resistance 2.}

\vspace{2mm}
\noindent\includegraphics[width=\linewidth]{\FIGDIR/f5_forest_2.png}
\figcap{The \textbf{3} difference-scale estimates: FEV$_1$ $-$36.8\,mL ($-$45.9 to $-$27.8) for PM$_1$ in
50,991 adults; HAPIN ITT HAZ $-$0.06 ($-$0.16 to 0.03, null); and the meta-analytic log-ratio of real-time
optical monitors to gravimetric mass, \textbf{0.31} (0.25--0.37), i.e.\ optical instruments read about
\textbf{36\,\%} high across 96 paired comparisons.}

\vspace{2mm}
\noindent\includegraphics[width=\linewidth]{\FIGDIR/f4_heatmap.png}
\figcap{Subtopic $\times$ architecture, pooled. The densest cells are \textit{Exposure} $\times$
\textit{Metadata only} and \textit{Exposure} $\times$ modelling; cohorts concentrate in \textit{Cardiovascular}
and \textit{Neuro}; sensing is split between instrument evaluation and field measurement.}

\vspace{4mm}

% -------------------------------------------------------------- venues + trials
\section*{\sffamily\bfseries\color{Deep}Venues, trial watch, and what the pipeline did}
\vspace{-1.5mm}{\color{Teal}\rule{\linewidth}{1.1pt}}\vspace{2.5mm}

\footnotesize
\textbf{Venues.} \textit{Environ Sci Technol} leads at \textbf{11} (9\,\%), then \textit{Atmos Environ} 8,
\textit{Atmosphere} 6, \textit{Research Square} 6, \textit{ACS Environ Health} 5, and \textit{Atmos Chem Phys}
and \textit{Atmos Pollut Res} at 4. \textbf{Preprints supplied 9} (Research Square 6, among them medRxiv with the HAPIN
follow-up, arXiv with Foehn).

\vspace{1.5mm}
\textbf{Recurring authors and most-cited work.} Surname frequency remains a naming artefact; \textbf{no
repeat-group cluster is claimed}. Citation counts for records days old are uninformative; the largest datasets
--- 65.2 and 26.3 million Medicare beneficiaries, 12.5 million admissions across six countries, 397,103 bacterial
isolates from 102 countries --- each supplied one record.

\vspace{1.5mm}
\textbf{Trial watch.} ClinicalTrials.gov \texttt{LastUpdatePostDate} windows 19--20, 20--21, 21--22 and 22--23
October: 0, 0, 2 and 2 hits. Stored: AIRWEIGHS NCT02763917 (HEPA purifier RCT, update) and the Taif COPD
prevalence survey NCT07886346 (observational, low relevance); APCIQ-BP NCT05903950 updated. \texttt{trials.json}
held \textbf{19 trials} over \textbf{35 windows} at period close. (A new controlled woodsmoke-exposure study,
NCT07536178, surfaced on the 30 October sweep and belongs to W44.) The week's interventional evidence again came
from the literature --- three crossover/RCT records --- not the registry.

\vspace{1.5mm}
\textbf{What the harvesting legs did.} \textbf{Europe PMC} lagged hard mid-week (3, 53, 1 records on 20--22
October) and then re-surfaced PubMed records one to two days late, producing many cross-day duplicates on
23--25 October, all caught at screening or by the DOI gate. \textbf{Crossref-by-ISSN} 7--69/day. \textbf{PubMed
connector} appended 2--20 records/day. \textbf{OpenAlex} HTTP 429 throughout. \textbf{arXiv} supplied Foehn,
held one day because its DataCite DOI was not yet registered. \textbf{Consensus} 10 hits, \textbf{0 new}; the DOI
gate caught one Consensus duplicate (Ginsburg 2026, shipped 29 July). \textbf{Scholar Gateway} not called.

\vspace{2mm}
% -------------------------------------------------------------- gaps
\section*{\sffamily\bfseries\color{Deep}Open gaps}
\vspace{-1.5mm}{\color{Teal}\rule{\linewidth}{1.1pt}}\vspace{2.5mm}

\footnotesize
\begin{itemize}
\item \textbf{The optical bias now has a pooled number and still no home.} Real-time optical monitors read
$\approx$36\,\% above gravimetric (log-ratio 0.31, 0.25--0.37); no health record in the week propagates that
bias, or calibration uncertainty, into an exposure--response estimate.

\item \textbf{Component effects rest on modelled components.} Every sulfate/nitrate/ammonium and source-specific
estimate this week comes from CTM, satellite--ML or receptor-model output; none from measured speciation at
the subject's location.

\item \textbf{Forecast skill is over-reported under random CV.} The GEOS-CF emulator's collapse under a
300\,km spatial block is a general warning for gridded PM products used to calibrate or gap-fill networks.

\item \textbf{The abstract gap is flat at 21\,\%}, concentrated in day-0 Elsevier deposits.

\item \textbf{Held queue empty.} The 3 October ES\&T NK-cell record entered the 23 October issue; Foehn was
released on 21 October.

\item \textbf{Pipeline defects carried forward}: scheduled runs for 23--29 October did not complete (recovered
on 30 October); Europe PMC late re-surfacing makes per-day dedup order-dependent; OpenAlex dead; Scholar Gateway
not reconnected; \texttt{.git} lock debris.
\end{itemize}

\clearpage

% -------------------------------------------------------------- digest
\section*{\sffamily\bfseries\color{Deep}Paper-level digest --- the week by subtopic}
\vspace{-1.5mm}{\color{Teal}\rule{\linewidth}{1.1pt}}\vspace{1mm}

{\fontsize{7.4}{9}\selectfont\color{Slate}Every record that appeared in a daily issue
keeps that issue's critical summary verbatim; nothing here is re-written or re-assessed.
Records surfaced without an abstract are shown as such. Ordered by subtopic, then by the
date of the issue that first carried the record.\par}

\input{digest_body}

% -------------------------------------------------------------- register
\section*{\sffamily\bfseries\color{Deep}Pooled corpus register}
\vspace{-1.5mm}{\color{Teal}\rule{\linewidth}{1.1pt}}\vspace{2.5mm}

\input{table_rows}

\vspace{4mm}
\begin{keybox}
{\sffamily\bfseries\footnotesize\color{Deep}Provenance and method}\\[1.2mm]
{\fontsize{7.2}{8.8}\selectfont
Period \textbf{18 -- 24 October 2026}, Sunday to Saturday, \textbf{seven daily issues},
\textbf{119 records}, \textbf{119 unique DOIs}, no record without a DOI. This rollup is
\textbf{aggregation only}: it re-reads \texttt{state/corpus/*.json}, \texttt{state/metrics.csv} and
\texttt{state/trials.json} over the period and \textbf{queries no literature API}; every record was gated by
\texttt{check\_dois.py} at daily-issue time. Tier mix \textbf{21 A, 43 B, 55 C}. The 23 and 24 October dailies and
this rollup were built on the \textbf{30 October recovery run} (the scheduled runs for 23--29 October did not
complete); no record is held. Effect estimates are transcribed verbatim from the daily issues on heterogeneous
contrasts and are \textbf{not} meta-analysed. Abstracts remain the copyright of their respective publishers.\par}
\end{keybox}

\end{document}
